%LaTeX format for the abstract
\documentclass[11pt]{article}
\usepackage{amsfonts,amsmath}

%\usepackage{amsmath}
\pagestyle{empty}
\textwidth=14 cm
\textheight=22 cm
\hoffset=-0.7in
\voffset=-0.6in
\parskip=6pt
\lineskip=18pt
%------------------------------------------------------------
\newcommand{\R}{\mathbb{R}}
\newcommand{\A}{\mathcal{A}}
\newcommand{\Q}{\mathcal{Q}}
\newcommand{\pt}{\partial_{t}}
\newcommand{\del}{\partial}
\newcommand{\pa}{\partial_{\alpha}}
\newcommand{\px}{\partial_{x}}
\newcommand{\e}{\epsilon}
\newcommand{\vh}{\widehat{v}}
\newcommand{\fh}{\widehat{F}}
\newcommand{\xh}{\widehat{\Xi}}
\newcommand{\tti}{\widetilde{T}}
\newcommand{\cof}{\mathop{\mathrm{cof}}}
\newcommand{\Div}{\mathop{\mathrm{Div}}}
\newcommand{\tr}{\mathop{\mathrm{tr}}}

\newcommand{\vb}{\bar{v}}
\newcommand{\fb}{\bar{F}}
\newcommand{\Sb}{\bar{S}}
\newcommand{\sh}{\widehat{S}}
\newcommand{\wh}{\widehat{W}}
\newcommand{\mue}{\frac{\mu}{\e}}
\newcommand{\intr}{\int_{\R^{3}}}

\begin{document}

Mathematical Methods in the Applied Sciences\\[0,3cm]


{\bf Subject:} submission to Mathematical Methods in the Applied Sciences\\[0,2cm]


Please find enclosed the pdf file of the article {\em``Weak solution of the merged mathematical equations of the polluted atmosphere''} by N. Juh\'asz and myself.
Considered as a geophysical fluid, the polluted atmosphere shares the shallow domain characteristics with other natural large-scale fluids such as seas and oceans. This means that its domain is excessively greater horizontally than in the vertical dimension, leading to the classic hydrostatic approximation of the Navier-Stokes equations. Essentially, the basic principle is to exploit the fact that when we are working on domains such as the ocean or the atmosphere, we are working on an �almost� two-dimensional set, and thus we can significantly simplify the model in the vertical direction, or even eliminate the third dimension and arrive to a limit model in two dimensions.  A common practice for example is to depth-integrate the Navier-Stokes equations for big lakes and seas and obtain the shallow water equations.
Our approach is different, we keep the vertical velocity and will arrive to a three dimensional limit model.
We use the
$$ \epsilon = \frac{\text{characteristic depth}}{\text{characteristic width}} $$
aspect ratio and the fact that, as the above described general ideas suggest, it is natural to consider asymptotic models and to take the limit model as $\epsilon$ goes to zero.
In this paper we consider the phenomena of a contaminant being emitted into, mixing with, and being deposited from the homogeneous and clean air. 

The purpose of this paper  is to create a self-contained model describing the atmospheric flow combined with the pollution effects and to investigate the asymptotic behaviour of the weak solutions of this model.  Specifically, as for the limit behavior of the weak solutions of the model we construct, we show that the weak solutions of this model converge to a weak solution of the hydrostatic limit model � in other words, taking the limit we arrive to a justification of the hydrostatic model for the case of the polluted atmosphere, where, instead of the originally used wind momentum equation, the hydrostatic air pressure assumption is used to describe the process in the vertical dimension.
The novelty of this present work is to provide a generalisation of previous  result (on the ocean) translated to the atmosphere, extending the fluid velocity equations with an additional convection-diffusion equation representing pollutants in the atmosphere.\\


   Thank you for you attention.
    
    Best regards,
    
    Donatella Donatelli
\end{document}









